\documentclass{exam}

\title{Circuit diagrams and boolean logic}
\author{}
\date{}

\usepackage{amsmath}
\usepackage{multicol}
\usepackage{circuitikz}
\usepackage{tikz}
\usepackage{float}
\usepackage{wrapfig}

\setlength{\columnsep}{1cm}

\usetikzlibrary{calc}

\newcommand{\answerbox}{%
  \begin{figure}[H]
    \begin{tikzpicture}
      \draw
      (0,0) rectangle (\textwidth, 0.45\textheight);
    \end{tikzpicture}
  \end{figure}
}

\tikzset{
  every path/.style={
    thick,
  },
  declare function={% in case of CVS which switches the arguments of atan2
    atan3(\a,\b)=ifthenelse(atan2(0,1)==90, atan2(\a,\b), atan2(\b,\a));
  },
  kinky cross radius/.initial=+.125cm,
  @kinky cross/.initial=+, kinky crosses/.is choice,
  kinky crosses/left/.style={@kinky cross=-},kinky crosses/right/.style={@kinky cross=+},
  kinky cross/.style args={(#1)--(#2)}{
    to path={
      let \p{@kc@}=($(\tikztotarget)-(\tikztostart)$),
          \n{@kc@}={atan3(\p{@kc@})+180} in
      -- ($(intersection of \tikztostart--{\tikztotarget} and #1--#2)!%
             \pgfkeysvalueof{/tikz/kinky cross radius}!(\tikztostart)$)
      arc [ radius     =\pgfkeysvalueof{/tikz/kinky cross radius},
            start angle=\n{@kc@},
            delta angle=\pgfkeysvalueof{/tikz/@kinky cross}180 ]
      -- (\tikztotarget)
    }
  }
}

\begin{document}

\maketitle

In this worksheet, you will practise your understanding of the relationship
between circuit diagrams and truth tables.

\begin{questions}
  \question
  Draw up the truth table for the following circuit diagrams.

  \begin{parts}
    \part
    Consider this diagram with three inputs and one output.

    \begin{figure}[H]
      \centering

      \begin{circuitikz}[x=3cm,y=1.5cm,z=2cm]
        \draw
        (0, 0)  node[not port] (not 1) {}
        (0, -1) node[not port] (not 2) {}
        (0, -2) node[not port] (not 3) {}
        ;
        \coordinate (bet 1 2) at ($ (not 1.out) !.5! (not 2.out) $);
        \coordinate (bet 2 3) at ($ (not 2.out) !.5! (not 3.out) $);
        \draw
        (not 1.in) ++(-1, 0) node (A) {A}
        (not 1.in) -- (A)
        (not 2.in) ++(-1, 0) node (B) {B}
        (not 2.in) -- (B)
        (not 3.in) ++(-1, 0) node (C) {C}
        (not 3.in) -- (C)
        ;
        \draw
        (bet 1 2)  ++(1, 0)  node[and port] (and 1 2) {}
        (bet 2 3)  ++(2, 0)  node[and port] (and 2 3) {}
        ;
        \draw
        (not 1.out)   |- (and 1 2.in 1)
        (not 2.out)   |- (and 1 2.in 2)
        (and 1 2.out) |- (and 2 3.in 1)
        (not 3.out)   |- (and 2 3.in 2)
        ;
        \draw
        (and 2 3.out) ++(1, 0) node (OUT) {OUT}
        (and 2 3.out) -- (OUT)
        ;
      \end{circuitikz}

      ~

      \begin{tabular}{|c|c|c||c|}
        \hline
        A & B & C & OUT \\ \hline
        ~ & ~ & ~ & ~ \\   \hline
        ~ & ~ & ~ & ~ \\   \hline
        ~ & ~ & ~ & ~ \\   \hline
        ~ & ~ & ~ & ~ \\   \hline
        ~ & ~ & ~ & ~ \\   \hline
        ~ & ~ & ~ & ~ \\   \hline
        ~ & ~ & ~ & ~ \\   \hline
        ~ & ~ & ~ & ~ \\   \hline
      \end{tabular}
    \end{figure}

    \vfill
    \part
    Consider this diagram with two inputs and one output.

    \begin{figure}[H]
      \centering

      \begin{circuitikz}[x=3cm,y=1.5cm,z=2cm]
        \draw
        (0, 0)      node           (A)       {A}
        %puts pen   a node         name of node    text at node
        (0, -1)     node           (B)       {B}
        (1, -0.5)   node[or port]  (orgate)  {}
                        %class for node, "or port" from package
        (1, -2)     node[and port] (and1)    {}
        (1.5, -2)   node[not port] (notgate) {}
        (2.5, -0.5) node[and port] (and2)    {}
        (3.0, -0.5) node           (OUTPUT)  {OUT}


        %node a, |- do these directions to destinatino orgate
        (A) |- (orgate.in 1)
        (B) |- (orgate.in 2)

        (orgate.out) |- (and2.in 1)
        (notgate.out) |- (and2.in 2)
        (and1.out) -- (notgate.in)
        (and2.out) -- (OUTPUT)
        ;
        \coordinate (foo) at ($ (orgate.in 1) - (0.05, 0) $);
        \draw
        (A) |- (foo)
            to[kinky cross=(B)--(orgate.in 2), kinky crosses=left]
            (and1.in 1 -| foo)
            -- (and1.in 1)
        (B) |- ($ (orgate.in 2) - (0.10, 0) $) |- (and1.in 2)
        ;
      \end{circuitikz}

      ~

      \begin{tabular}{|c|c||c|}
        \hline
        A & B & OUT \\ \hline
        ~ & ~ & ~ \\ \hline
        ~ & ~ & ~ \\ \hline
        ~ & ~ & ~ \\ \hline
        ~ & ~ & ~ \\ \hline
      \end{tabular}
    \end{figure}

    \part
    Consider this circuit with two inputs and two outputs.

    \begin{figure}[H]
      \centering

      \begin{circuitikz}[x=3cm,y=1.5cm,z=2cm]
        \draw
        (0, 0)  node [xor port] (xor) {}
        (0, -1) node [and port] (and) {}
        ;
        \draw
        (xor.out) ++(1, 0) node (S) {S}
        (xor.out) -- (S)
        (and.out) ++(1, 0) node (C) {C}
        (and.out) -- (C)
        ;
        \draw
        (xor.in 1) ++(-1, 0) node (A) {A}
        (xor.in 2) ++(-1, 0) node (B) {B}
        (xor.in 1) -- (A)
        (xor.in 2) -- (B)
        ;
        \draw
        (xor.in 1) ++(-0.1, 0) coordinate (foo)
          to[kinky cross=(A)--(xor.in 2), kinky crosses=left]
          (and.in 1 -| foo)
          -- (and.in 1)
        (xor.in 2) ++(-0.2, 0)
          |- (and.in 2)
        ;
      \end{circuitikz}

      ~

      \begin{tabular}{|c|c||c|c|}
        \hline
        A & B & S & C \\ \hline
        ~ & ~ & ~ & ~ \\ \hline
        ~ & ~ & ~ & ~ \\ \hline
        ~ & ~ & ~ & ~ \\ \hline
        ~ & ~ & ~ & ~ \\ \hline
      \end{tabular}
    \end{figure}
  \end{parts}

  \pagebreak
  \question
  Draw a circuit diagram that implements the boolean function.

  \begin{parts}
    \part
    Consider the following function with two inputs and one output.

    \begin{figure}[H]
      \centering
      \begin{tabular}{cc||c}
        A & B & OUT \\ \hline
        $1$ & $1$ & $0$ \\
        $1$ & $0$ & $0$ \\
        $0$ & $1$ & $0$ \\
        $0$ & $0$ & $1$ \\
      \end{tabular}
    \end{figure}

    \answerbox

    \vfill
    \part
    Consider the following function with three inputs and one output.

    \begin{figure}[H]
      \centering
      \begin{tabular}{ccc||c}
        A & B & C & OUT \\ \hline
        $0$ & $0$ & $0$ & $0$ \\
        $0$ & $0$ & $1$ & $1$ \\
        $0$ & $1$ & $0$ & $1$ \\
        $0$ & $1$ & $1$ & $1$ \\
        $1$ & $0$ & $0$ & $1$ \\
        $1$ & $0$ & $1$ & $1$ \\
        $1$ & $1$ & $0$ & $1$ \\
        $1$ & $1$ & $1$ & $1$ \\
      \end{tabular}
    \end{figure}

    \answerbox
  \end{parts}
\end{questions}

\end{document}
